% =====================================================================
%  Module 4 : Stockage, bases de données et messages
% =====================================================================
\section{Stockage, bases de données et messages}

\begin{frame}{Trois types de stockage}
  \centering
  \includegraphics[width=.96\textwidth]{stockage-types}
\end{frame}

\begin{frame}{Amazon S3 : Simple Storage Service}
  \begin{itemize}
    \setlength{\itemsep}{.45em}
    \item Stockage objet, accessible par API HTTPS
    \item \textbf{Objet} : clé + données + métadonnées ; de 0 octet à 50 To
    \item Envoi en une requête jusqu'à 5 Go, au-delà téléversement en plusieurs parties
    \item \textbf{Bucket} : conteneur d'objets, créé dans une région
    \item Nom de bucket unique au monde (3 à 63 caractères, minuscules)
    \item Cohérence forte en lecture après écriture depuis 2020
    \item Plus de 500 000 milliards d'objets stockés, plus de 200 millions de requêtes par seconde
  \end{itemize}
  \source{AWS, \emph{Amazon S3 FAQs} ; S. Stormacq, « Twenty years of Amazon S3 », AWS News Blog, 2026.}
\end{frame}

\begin{frame}{Buckets, clés et préfixes}
  \centering
  \includegraphics[width=.96\textwidth]{s3-modele}\par\medskip
  \raggedright
  \begin{itemize}
    \item Pas de dossiers : des clés à plat, regroupées par préfixe dans la console
  \end{itemize}
\end{frame}

\begin{frame}{Classes de stockage S3}
  \renewcommand{\arraystretch}{1.12}
  \begin{tabularx}{\textwidth}{@{}>{\raggedright\arraybackslash}p{5.4cm}>{\raggedright\arraybackslash}X@{}}
    \toprule
    \textbf{Classe} & \textbf{Usage} \\ \midrule
    S3 Standard & données lues fréquemment (classe par défaut) \\
    S3 Intelligent-Tiering & accès imprévisible, niveaux automatiques \\
    S3 Standard-IA & lectures rares, accès immédiat \\
    S3 One Zone-IA & comme Standard-IA, dans une seule AZ \\
    S3 Glacier Instant Retrieval & archives lues environ une fois par trimestre \\
    S3 Glacier Flexible Retrieval & archives, récupération en minutes à heures \\
    S3 Glacier Deep Archive & archives longues, récupération en 12 à 48 h \\
    S3 Express One Zone & latence très faible, une seule AZ \\ \bottomrule
  \end{tabularx}
  \source{\url{https://aws.amazon.com/s3/storage-classes/}}
\end{frame}

\begin{frame}{Durabilité et disponibilité}
  \begin{itemize}
    \setlength{\itemsep}{.55em}
    \item Durabilité conçue pour 99,999999999 \% par an (11 neuf)
    \item Soit une probabilité de perte de \(10^{-11}\) par objet et par an
    \item Avec 10 millions d'objets : une perte attendue tous les 10 000 ans environ
    \item Données copiées sur plusieurs AZ (sauf classes One Zone)
    \item SLA de disponibilité de S3 Standard : 99,9 \%
  \end{itemize}
  \bigskip
  \attention La durabilité protège contre les pannes matérielles, pas contre une suppression : activer le versionnage
\end{frame}

\begin{frame}{Fonctionnalités de S3}
  \begin{itemize}
    \setlength{\itemsep}{.45em}
    \item \textbf{Versionnage} : conservation des versions précédentes de chaque objet
    \item \textbf{Cycle de vie} : changement de classe ou suppression automatique selon l'âge
    \item \textbf{Réplication} : vers un autre bucket, dans la même région ou une autre
    \item \textbf{Chiffrement} : SSE-S3 appliqué par défaut depuis janvier 2023, SSE-KMS en option
    \item \textbf{Notifications d'événements} : vers SQS, SNS, Lambda, EventBridge
    \item \textbf{Hébergement de site statique}
    \item \textbf{Object Lock} : écriture unique, lecture multiple (WORM)
  \end{itemize}
\end{frame}

\begin{frame}{Contrôle d'accès à S3}
  \begin{itemize}
    \setlength{\itemsep}{.45em}
    \item Stratégies IAM de l'appelant (utilisateur, rôle)
    \item Stratégie de bucket (basée sur la ressource)
    \item \textbf{Block Public Access} : activé par défaut sur les nouveaux buckets depuis avril 2023
    \item ACL désactivées par défaut (\emph{Object Ownership : bucket owner enforced})
    \item URL présignée : accès temporaire à un objet précis
  \end{itemize}
  \bigskip
  \attention Juin 2017 : bucket public d'un sous-traitant de Verizon, données de 6 à 14 millions de clients exposées
  \source{AWS, « Amazon S3 now applies two security best practices to all new buckets by default », 2023 ; UpGuard, « Cloud Leak: How A Verizon Partner Exposed Millions of Customer Accounts », 2017.}
\end{frame}

\begin{frame}{Qui peut lire un objet ?}
  \centering
  \includegraphics[width=.95\textwidth]{s3-acces}
\end{frame}

\begin{frame}[fragile]{Stratégie de bucket}
  Exemple : refuser toute requête qui n'utilise pas HTTPS
\begin{lstlisting}[language=json]
{
  "Version": "2012-10-17",
  "Statement": [{
    "Sid": "HTTPSUniquement",
    "Effect": "Deny",
    "Principal": "*",
    "Action": "s3:*",
    "Resource": ["arn:aws:s3:::galerie-camille",
                 "arn:aws:s3:::galerie-camille/*"],
    "Condition": { "Bool": { "aws:SecureTransport": "false" } }
  }]
}
\end{lstlisting}
  \lien{https://docs.aws.amazon.com/AmazonS3/latest/userguide/example-bucket-policies.html}
\end{frame}

\begin{frame}{Tarification de S3}
  \begin{itemize}
    \setlength{\itemsep}{.5em}
    \item Volume stocké par mois, selon la classe : un peu plus de 0,02 \$/Go-mois en Standard
    \item Requêtes : \texttt{PUT}, \texttt{COPY}, \texttt{LIST} plus chères que \texttt{GET}
    \item Transfert sortant vers Internet ; transfert entrant gratuit
    \item Frais de récupération pour les classes IA et Glacier
    \item Fonctionnalités optionnelles : réplication, analyse, Transfer Acceleration
  \end{itemize}
  \lien{https://aws.amazon.com/s3/pricing/}
\end{frame}

\begin{frame}{Amazon EBS : Elastic Block Store}
  \begin{itemize}
    \setlength{\itemsep}{.45em}
    \item Volume de stockage en bloc, vu comme un disque par l'instance
    \item Dans une seule AZ, répliqué à l'intérieur de cette AZ
    \item Attaché à une instance à la fois (sauf Multi-Attach pour \texttt{io1}/\texttt{io2})
    \item Système de fichiers à créer par l'utilisateur
    \item Instantanés incrémentaux stockés dans S3, copiables entre régions
    \item Chiffrement des volumes et des instantanés (KMS)
  \end{itemize}
  \bigskip
  \attention Pas de multi-AZ : un instantané permet de recréer le volume dans une autre AZ
\end{frame}

\begin{frame}{Types de volumes EBS}
  \renewcommand{\arraystretch}{1.12}
  \begin{tabularx}{\textwidth}{@{}>{\raggedright\arraybackslash}p{3.6cm}>{\raggedright\arraybackslash}p{1.6cm}>{\raggedright\arraybackslash}X >{\raggedright\arraybackslash}p{2.6cm}@{}}
    \toprule
    \textbf{Type} & \textbf{Support} & \textbf{Usage} & \textbf{Taille max.} \\ \midrule
    \texttt{gp3} & SSD & usage général, 3 000 à 80 000 IOPS & 64 Tio \\
    \texttt{gp2} & SSD & génération précédente & 16 Tio \\
    \texttt{io2} Block Express & SSD & bases critiques, jusqu'à 256 000 IOPS & 64 Tio \\
    \texttt{st1} & HDD & débit séquentiel (journaux) & 16 Tio \\
    \texttt{sc1} & HDD & données froides, coût minimal & 16 Tio \\ \bottomrule
  \end{tabularx}
  \source{AWS, \emph{Amazon EBS volume types} ; « Amazon EBS increases the maximum size and provisioned performance of gp3 volumes », 2025.}
\end{frame}

\begin{frame}{Amazon EFS et comparaison}
  \begin{itemize}
    \item EFS : système de fichiers NFS managé, élastique, partagé par plusieurs instances, sur plusieurs AZ
  \end{itemize}
  \medskip
  \renewcommand{\arraystretch}{1.12}
  \begin{tabularx}{\textwidth}{@{}>{\raggedright\arraybackslash}p{2.6cm}>{\raggedright\arraybackslash}X>{\raggedright\arraybackslash}X>{\raggedright\arraybackslash}X@{}}
    \toprule
    & \textbf{S3} & \textbf{EBS} & \textbf{EFS} \\ \midrule
    Type & objet & bloc & fichiers (NFS) \\
    Accès & API HTTPS & une instance & plusieurs instances \\
    Redondance & plusieurs AZ (selon classe) & une AZ & plusieurs AZ \\
    Facturation & volume réel + requêtes & volume provisionné & volume réel \\
    Usage type & fichiers, sauvegardes & disque système, base & dossier partagé \\ \bottomrule
  \end{tabularx}
\end{frame}

% ---- bases de données ------------------------------------------------------
\begin{frame}{Bases relationnelles et NoSQL}
  \renewcommand{\arraystretch}{1.35}
  \begin{tabularx}{\textwidth}{@{}>{\raggedright\arraybackslash}p{2.8cm}>{\raggedright\arraybackslash}X>{\raggedright\arraybackslash}X@{}}
    \toprule
    & \textbf{Relationnel (SQL)} & \textbf{NoSQL} \\ \midrule
    Modèle & tables, schéma fixe, jointures & clé-valeur, document, graphe ; schéma souple \\
    Requêtes & SQL & API propre à chaque moteur \\
    Mise à l'échelle & surtout verticale & horizontale (partitionnement) \\
    Garanties & transactions ACID & souvent cohérence à terme, transactions limitées \\
    Exemples & PostgreSQL, MySQL, Oracle & DynamoDB, MongoDB, Redis, Neo4j \\ \bottomrule
  \end{tabularx}
\end{frame}

\begin{frame}{Bases de données managées sur AWS}
  \renewcommand{\arraystretch}{1.35}
  \begin{tabularx}{\textwidth}{@{}>{\raggedright\arraybackslash}p{4cm}>{\raggedright\arraybackslash}X@{}}
    \toprule
    \textbf{Amazon RDS} & PostgreSQL, MySQL, MariaDB, Oracle, SQL Server, Db2 \\
    \textbf{Amazon Aurora} & compatible PostgreSQL et MySQL, stockage réparti sur 3 AZ \\
    \textbf{Amazon DynamoDB} & clé-valeur et document, serverless \\
    \textbf{Amazon ElastiCache} & cache en mémoire : Valkey, Redis OSS, Memcached \\
    \textbf{Amazon DocumentDB} & document, compatible MongoDB \\
    \textbf{Amazon Neptune} & graphe \\
    \textbf{Amazon OpenSearch Service} & recherche et analyse de journaux \\ \bottomrule
  \end{tabularx}
  \lien{https://aws.amazon.com/products/databases/}
\end{frame}

\begin{frame}{Amazon RDS}
  \begin{columns}[T]
    \begin{column}{.48\textwidth}
      \textbf{Géré par AWS}
      \begin{itemize}
        \setlength{\itemsep}{.35em}
        \item installation et correctifs du moteur
        \item sauvegardes automatiques (stockées dans S3), restauration à un instant donné
        \item copie de secours synchrone (Multi-AZ)
        \item chiffrement, supervision
      \end{itemize}
    \end{column}
    \begin{column}{.48\textwidth}
      \textbf{Choisi par le client}
      \begin{itemize}
        \setlength{\itemsep}{.35em}
        \item moteur et version
        \item classe d'instance (\texttt{db.t4g.micro}, \texttt{db.r7g.large}...)
        \item stockage (gp3, io2)
        \item Multi-AZ, réplicas en lecture
        \item réseau, Security Group, utilisateurs
      \end{itemize}
    \end{column}
  \end{columns}
  \bigskip
  \attention Facturée à l'heure tant qu'elle existe, même sans requête
\end{frame}

\begin{frame}{Amazon DynamoDB}
  \begin{itemize}
    \setlength{\itemsep}{.45em}
    \item Base clé-valeur et document, entièrement managée, sans serveur à dimensionner
    \item Latence de quelques millisecondes quelle que soit la taille de la table
    \item Clé primaire : clé de partition, éventuellement clé de tri
    \item Élément de 400 Ko au maximum
    \item Facturation à la requête (\emph{on-demand}) ou capacité réservée
    \item Index secondaires, transactions, flux de modifications (Streams)
  \end{itemize}
  \bigskip
  \attention Conception guidée par les requêtes à servir : une requête imprévue est difficile à réaliser
  \lien{https://docs.aws.amazon.com/amazondynamodb/latest/developerguide/Introduction.html}
\end{frame}

\begin{frame}{S3, RDS ou DynamoDB ?}
  \begin{columns}[c]
    \begin{column}{.55\textwidth}
      \includegraphics[width=\textwidth]{galerie-evoluee}
    \end{column}
    \begin{column}{.43\textwidth}
      \begin{itemize}
        \setlength{\itemsep}{.5em}
        \item \textbf{S3} : fichiers (images, documents, sauvegardes)
        \item \textbf{RDS} : données structurées, jointures, transactions
        \item \textbf{DynamoDB} : accès par clé à grande échelle, schéma souple
        \item Critères : volume, fréquence d'accès, type de requêtes, cohérence, coût
      \end{itemize}
    \end{column}
  \end{columns}
\end{frame}

% ---- messages -----------------------------------------------------------------
\begin{frame}{Amazon SQS : Simple Queue Service}
  \begin{itemize}
    \setlength{\itemsep}{.45em}
    \item File de messages managée, lancée en 2006
    \item Découple un \textbf{producteur} et un \textbf{consommateur}
    \item Absorbe les pics de charge, tolère l'arrêt d'un consommateur
    \item Opérations : \texttt{SendMessage}, \texttt{ReceiveMessage}, \texttt{DeleteMessage}
    \item Message jusqu'à 1 Mio (256 Kio avant août 2025)
    \item Conservation : 4 jours par défaut, de 1 minute à 14 jours
  \end{itemize}
  \lien{https://docs.aws.amazon.com/AWSSimpleQueueService/latest/SQSDeveloperGuide/welcome.html}
\end{frame}

\begin{frame}{Délai de visibilité}
  \centering
  \includegraphics[width=.96\textwidth]{sqs-visibilite}\par\medskip
  \raggedright
  \begin{itemize}
    \item Message reçu = invisible pendant le délai (30 s par défaut, jusqu'à 12 h)
    \item Supprimé seulement par \texttt{DeleteMessage} ; sinon il réapparaît
    \item Conséquence : traitement idempotent côté consommateur
  \end{itemize}
\end{frame}

\begin{frame}{Files standard et FIFO}
  \renewcommand{\arraystretch}{1.4}
  \begin{tabularx}{\textwidth}{@{}>{\raggedright\arraybackslash}p{3.2cm}>{\raggedright\arraybackslash}X>{\raggedright\arraybackslash}X@{}}
    \toprule
    & \textbf{Standard} & \textbf{FIFO} \\ \midrule
    Livraison & au moins une fois & traitement unique \\
    Ordre & au mieux & strict (par groupe de messages) \\
    Débit & quasi illimité & 300 messages/s par action, 3 000 par lots ; plus en mode haut débit \\
    Nom & libre & suffixe \texttt{.fifo} \\ \bottomrule
  \end{tabularx}
  \bigskip
  \begin{itemize}
    \item File de lettres mortes (DLQ) : messages en échec après N réceptions
    \item Attente longue (\emph{long polling}) jusqu'à 20 s : moins de requêtes vides
  \end{itemize}
\end{frame}

\begin{frame}{SQS, SNS, EventBridge}
  \renewcommand{\arraystretch}{1.4}
  \begin{tabularx}{\textwidth}{@{}>{\raggedright\arraybackslash}p{3.2cm}>{\raggedright\arraybackslash}X>{\raggedright\arraybackslash}X@{}}
    \toprule
    \textbf{Service} & \textbf{Modèle} & \textbf{Usage} \\ \midrule
    SQS & file : un message, un consommateur & tâches en arrière-plan, lissage de charge \\
    SNS & publication-abonnement : un message, plusieurs abonnés & notifications, diffusion vers plusieurs files \\
    EventBridge & bus d'événements avec règles de routage & intégration entre services et applications \\ \bottomrule
  \end{tabularx}
  \bigskip
  \begin{itemize}
    \item SQS : premier million de requêtes gratuit chaque mois, puis environ 0,40 \$ par million (standard)
  \end{itemize}
  \source{AWS, \emph{Amazon SQS pricing}.}
\end{frame}

\diapotp{4}{Un bucket, un rôle, une file}{%
  \item Bucket privé, téléversement depuis la console
  \item URL présignée
  \item Instance EC2 avec rôle IAM : lecture et écriture dans S3
  \item File SQS : envoi, réception, délai de visibilité}

